\chapter{Second Model Family}
\label{app:second}

To test whether the results are specific to Qwen3-8B, a second family,
DeepSeek-R1-Distill-Llama-8B (8B), was fine-tuned with the same data, LoRA
configuration, and evaluation protocol (one seed). The model was trained to emit
an empty thinking block followed by the score, because the base model otherwise
produces long rationales and does not answer within a short generation budget.

\begin{table}[h]
\centering
\small
\caption{Second-family comparison. In-distribution rows are 2016--2019; 2024 is
the temporal test. Choice chance is 0.163 in-distribution and 0.500 in 2024.}
\label{tab:second}
\begin{tabular}{lcccccc}
\toprule
& \multicolumn{3}{c}{2016--2019} & \multicolumn{3}{c}{2024} \\
\cmidrule(lr){2-4}\cmidrule(lr){5-7}
Model & MAE & corr & choice & MAE & corr & choice \\
\midrule
Qwen3-8B (3 seeds) & 20.84 & 0.679 & 0.518 & 36.47 & $-0.029$ & 0.527 \\
DeepSeek-R1-Distill (1 seed) & 22.66 & 0.635 & 0.548 & 41.44 & $-0.017$ & 0.523 \\
HistGB & 20.58 & 0.714 & 0.534 & 35.06 & $-0.030$ & 0.456 \\
\bottomrule
\end{tabular}
\end{table}

The qualitative pattern is the same for both families: there is real
in-distribution signal (correlation 0.635 for DeepSeek versus 0.679 for Qwen),
and the 2024 correlation collapses to near zero ($-0.017$ versus $-0.029$). The
Qwen model has a lower MAE and a higher correlation; the DeepSeek model has a
slightly higher choice accuracy. With one seed for the second family, these
differences should be read as descriptive. The result supports the claim that the
in-distribution signal and the temporal failure are not specific to a single
model family, while a broader multi-family study remains future work.
